\section{Анализ органов РНМ других \commentnote{DISS-005-C0033}{Евдаков Алексей}{Наверное стоит фирмы перенести в меньший уровень, или как-то иначе описать, ибо сейчас это заголовки, что слишком «глобально».}{производителей }}
\label{sec:ch1_manufacturers_analysis}

В данном разделе приводятся формулы или описания РНМ у различных производителей и оценка их поведения с точки зрения описанных выше базовых и современных проблем функционирования.

\subsection{\commentnote{DISS-005-C0034}{Евдаков Алексей}{«\_вспомогательные файлы\textbackslash{}\_РЭ\textbackslash{}Экра», Файл «РЭ\_БАВР» страница 18 (есть пропущенные)}{ЭКРА}}

\subsubsection{Описание органа}

Выполнен в фазном исполнении по схеме -- фаз $I_A$ и $U_{BC}$. Задаётся через максимальные и минимальные углы, а угол максимальной чувствительности является их биссектрисой. Амплитуда не указывается, так как это отдельный орган, выполняющий простейшую функцию сравнения с уставкой.

\textbf{\commentnote{DISS-005-C0035}{Евдаков Алексей}{Может не красиво, но реально, нафига скрывать РУКОВОДСТВА? Это же базовые документы.}{Примечание}:} какая-либо информация скрывается данным производителем от общественности. Руководство эксплуатации на данные устройства невозможно найти в открытом доступе. Поэтому полноценный научным сообществом сложен. И во многом оценка приводится на основании анализа реальных осциллограмм функционирования.

\begin{figure}[htbp]
    \centering
    \includegraphics[width=0.7\textwidth]{image17.png}
    \caption[РНМ производства ЭКРА]{РНМ производства ЭКРА.}
    \label{fig:ch1_img17}
\end{figure}

\subsubsection{Оценка функционирования}

Данная реализация имеет только фазное исполнение, что, согласно современным тенденциям, является не лучшим решением для сетей с изолированной нейтралью и может приводить к излишним срабатывания.

Характеристика срабатывания по своей схеме близка к базовой в первых версиях исполнения с тем отличием, что ширина зоны срабатывания задаётся не фиксированной -- 180 градусов, а может уменьшаться. Данная особенность полезная, однако не решает концептуально ни одной из проблем, и всё равно требует ручной подстройки. Логика функционирования скопирована с стандартных органов РЗА без существенных модернизаций.

Ни базовые, ни современные проблемы в РНМ данного производителя не решены.

\subsection{Механотроника (\commentnote{DISS-005-C0036}{Евдаков Алексей}{«\_вспомогательные файлы\textbackslash{}\_РЭ\textbackslash{}БМРЗ»}{БМРЗ})}

\subsubsection{Описание органа}

Выполнен в фазном исполнении по схеме -- фаз $I_A$ и $U_{BC}$. Задаётся через угол максимальной чувствительности и отстройкой от него на $\pm 90$ градусов. Присутствует угол гистерезис по углу. При потере питания используется напряжение памяти -- 100~мс.

\begin{figure}[htbp]
    \centering
    \includegraphics[width=0.7\textwidth]{image18.png}
    \caption[РНМ производства Механотроника]{РНМ производства Механотроника.}
    \label{fig:ch1_img18}
\end{figure}

\subsubsection{Оценка функционирования}

Данная реализация имеет только фазное исполнение, что, согласно современным тенденциям, является не лучшим решением для сетей с изолированной нейтралью и может приводить к излишним срабатывания.

Характеристика срабатывания по своей схеме близка к базовой в первых версиях исполнения с тем отличием, что присутствует гистерезис на границе, который изменяет поведение. Данная особенность полезная, однако не решает концептуально ни одной из проблем, и всё равно требует ручной подстройки. А также используется не напряжение соседней секции, а напряжение памяти, что может приводить к проблемам функционирования при длительной аварии.

Ни базовые, ни современные проблемы в РНМ данного производителя не решены.

\subsection{\commentnote{DISS-005-C0037}{Евдаков Алексей}{взято с их сайта\newline «\_вспомогательные файлы\textbackslash{}\_РЭ\textbackslash{}Бреслер»}{Бреслер}}

\subsubsection{Описание органа}

Выполнен на основе анализа прямых последовательностей. Характеристики срабатывания описаны в \commentnote{DISS-005-C0038}{Евдаков Алексей}{Согласно локальным файлам, аналогичен нашему. Однако вордовский файл может быть конфиденциальным. И в широком доступе не нашёл.\newline НО есть в диссертации Жукова, так что норм.\newline Файл хороший, есть подробная математика.}{формуле }ХХ взятой из диссертации Жукова «Повышение эффективности работы АВР» [ХХ\commentnote{DISS-005-C0039}{Евдаков Алексей}{Вставить ссылку}{]}]. Приведённые характеристики в точности совпадают, поэтому решено приравнять их. В качестве уставок задаются только: уставка по току, коэффициент подпитки, угол максимальной чувствительности.

,

\begin{figure}[htbp]
    \centering
    \includegraphics[width=0.7\textwidth]{image20.png}
    \label{fig:ch1_img20}
\end{figure}

\subsubsection{Оценка функционирования}

Данная реализация имеет только исполнение только по прямой последовательности, что, согласно современным тенденциям актуально для сетей с изолированной нейтралью (6-35кВ). Однако применение в сетях с глухозаземлённой нейтралью (0,4кВ) может вызвать сложности. Характеристика срабатывания по своей схеме практически аналогична к базовой в первых версиях исполнения. Имеется интересное решение с использование активной составляющей полученной характеристики, однако данный фактор не вносит существенных отличий по отношению к заданию двух угловых порогов или угла максимальной чувствительности

Ни базовые, ни современные проблемы в РНМ данного производителя не решены.
